\section{The delay of a contact: who follows whom}\label{sec:lag}

The groups of Section~\ref{sec:groups} are built from the contact alone. To say who follows whom
(Sections~\ref{sec:leaders} and~\ref{sec:comparison}) we also need how many seconds $j$ is behind
$i$. As in~\cite{li2013}, at each instant $t$ with $m_{ij}(t)=1$ we take the delay at which $j$
passes closest to where $i$ is,
\begin{equation}\label{eq:lstar}
  l^*(t)=\arg\min_{l:\,(t,t+l)\in\M_{ij}}\|\x_i(t)-\x_j(t+l)\| ,
\end{equation}
and over a pair interval $[t_{\rm on},t_{\rm off})$ (Definition~\ref{def:pair-interval}) the
signed median delay
\[
  \ell_{\rm med}=\mathrm{median}\{l^*(t):\ t\in[t_{\rm on},t_{\rm off}),\ m_{ij}(t)=1\},
\]
with $\ell_{\rm med}>0$ when $j$ follows $i$. How well $l^*$ is defined depends on the phase of
flight.

\paragraph{Glide.} The path is nearly straight and the distance has a single sharp minimum: $l^*$
is the delay at which $j$ crosses the point of $i$, $+120$~s in Figure~\ref{fig:contact}.

\paragraph{Thermal.} $j$ circles over one spot, so the distance $\|\x_i(t)-\x_j(t+l)\|$, as a
function of $l$, has a nearly equal minimum at every turn (Figure~\ref{fig:clock}, top right): the
minimiser is degenerate, and $l^*(t)$ falls anywhere among the accepted delays. Altitude narrows
them. In a climb at rate $w$, pilot $j$ gains $wl$ metres in $l$ seconds,
\begin{equation}\label{eq:climb}
  z_j(t+l)\simeq z_j(t)+w\,l ,
\end{equation}
so, with $\Delta z = z_i(t)-z_j(t)$ the height of $i$ above $j$, the vertical condition
$|z_i(t)-z_j(t+l)|<\varepsilon_z$ becomes $|\Delta z-w\,l|<\varepsilon_z$, that is
\begin{equation}\label{eq:window}
  l\in\left(\frac{\Delta z-\varepsilon_z}{w},\ \frac{\Delta z+\varepsilon_z}{w}\right):
\end{equation}
an interval of width $2\varepsilon_z/w$ around $\Delta z/w$, the time $j$ needs to reach the
altitude of $i$, which the scan cuts to $[-\delta,\delta]$. With the starting $\delta=5$~min this
at least halves the range of $l^*$. With the chosen $\delta=2$~min the interval, 300~s wide at
$w=2$~m/s, is wider than $[-\delta,\delta]$, and altitude gains little: at $\Delta z=200$~m, from
$[-120,120]$ to $(-50,120]$~s.

\paragraph{Search.} The pilot turns without climbing, $w\simeq0$, and altitude constrains no delay.
The instants at which either pilot is in \st{search} are left out of the median; $\ell_{\rm med}$
is undefined when none remains.

\begin{figure}
\centering
\animated{altitude_clock}
\caption{\animmark The delay in a thermal, for one instant $t$ and every delay $l$; the animation
runs the delay $l$ from $-300$ to $300$~s. Left: $i$ and $j$ in the same thermal, $j$ lower by
$\Delta z=200$~m; the band is $z_i(t)\pm\varepsilon_z$. Top right: the horizontal distance between
$\x_i(t)$ and $\x_j(t+l)$; inside the thermal it never exceeds 50~m, far below
$\varepsilon_h=500$~m, so the horizontal condition alone accepts every delay of the stay; it has
an equal minimum at every turn, so it cannot fix $l^*$. Middle right: the altitude of $j$ at $t+l$
crosses the band in a single window. Bottom: the delays that satisfy each condition and both; both
hold for $-50<l<250$~s, the interval~\eqref{eq:window}, centred on $\Delta z/w=100$~s, and the scan
keeps those with $|l|\le\delta$: $-50<l\le120$~s with $\delta=2$~min.
Before $l=-120$~s, $j$ was gliding towards the thermal: at $l\approx-280$~s it had the right
altitude but was 1.4~km away. Model: $j$ circles with radius 25~m, one turn every 20~s, climbing at
$w=2$~m/s; it arrives at 9~m/s, sinking at 1~m/s; $i$ is on the same circle.}
\label{fig:clock}
\end{figure}
